\subsection{Polynomial Spectral Mapping}

\begin{proposition} \label{prop:eigenvalues-of-polynomial-in-matrix} \hlt{Eigenvalues of a Polynomial in a Matrix}\\
Let $A\in\operatorname{Mat}(n,n,\mathbb{C})$, let $p\in\mathbb{C}[z]$, and let $\lambda_1,\ldots,\lambda_n$ be the eigenvalues of $A$, counted with algebraic multiplicity. Then the eigenvalues of $p(A)$ are $p(\lambda_1),\ldots,p(\lambda_n)$, counted with algebraic multiplicity.
\begin{proof}
Choose a Schur form $U^*AU=T$ with diagonal $\lambda_1,\ldots,\lambda_n$. Then $p(A)=Up(T)U^*$, and $p(T)$ is upper triangular with diagonal $p(\lambda_1),\ldots,p(\lambda_n)$. Hence $p(A)$ has precisely these eigenvalues.\\
\hlt{Structural Commentary}: Schur form reduces the claim to the diagonal.\\
\hlt{Key Insight}: A polynomial preserves triangular form and applies $p$ to each diagonal entry.
\end{proof}
\end{proposition}

\begin{corollary} \label{cor:polynomial-spectral-mapping-matrix} \hlt{Polynomial Spectral Mapping}\\
Let $A\in\operatorname{Mat}(n,n,\mathbb{C})$ and $p\in\mathbb{C}[z]$. Then $\operatorname{spec}(p(A))=\{p(\lambda):\lambda\in\operatorname{spec}(A)\}$.
\begin{proof}
Let $\lambda_1,\ldots,\lambda_n$ be the eigenvalues of $A$, counted with algebraic multiplicity.\\
By Proposition \ref{prop:eigenvalues-of-polynomial-in-matrix}, the eigenvalues of $p(A)$, counted with algebraic multiplicity, are $p(\lambda_1),\ldots,p(\lambda_n)$. Taking underlying sets gives the claim.\\
\hlt{Structural Commentary}: Polynomial evaluation pushes the spectrum forward through the scalar polynomial.\\
\hlt{Key Insight}: Once Schur form has exposed the diagonal, $p$ acts entry by entry there.
\end{proof}
\end{corollary}

\begin{corollary} \label{cor:polynomial-spectral-mapping-operator} \hlt{Polynomial Spectral Mapping for Operators}\\
Let $\mathcal{V}\neq\{\mathbf{0}\}$ be finite-dimensional over $\mathbb{C}$, let $T\in\LL(\mathcal{V})$, and let $p\in\mathbb{C}[z]$.\\
Then $\operatorname{spec}(p(T))=\{p(\lambda):\lambda\in\operatorname{spec}(T)\}$.
\begin{proof}
Choose a basis $\beta$ and set $A=[T]_{\beta\leftarrow\beta}$.\\
Then $[p(T)]_{\beta\leftarrow\beta}=p(A)$, and coordinate isomorphism gives $\operatorname{spec}(T)=\operatorname{spec}(A)$ and $\operatorname{spec}(p(T))=\operatorname{spec}(p(A))$.\\
By Corollary \ref{cor:polynomial-spectral-mapping-matrix}, $\operatorname{spec}(p(T))=\operatorname{spec}(p(A))=\{p(\lambda):\lambda\in\operatorname{spec}(A)\}=\{p(\lambda):\lambda\in\operatorname{spec}(T)\}$.\\
\hlt{Structural Commentary}: Coordinates reduce the operator statement to the matrix statement.\\
\hlt{Key Insight}: Polynomial spectral mapping is an operator fact seen through any convenient basis.
\end{proof}
\end{corollary}

\begin{corollary} \label{cor:polynomial-in-matrix-invertible-spectrum-test} \hlt{Polynomial Invertibility Test}\\
Let $A\in\operatorname{Mat}(n,n,\mathbb{C})$ and $p\in\mathbb{C}[z]$. If $p(\lambda)\neq0$ for every $\lambda\in\operatorname{spec}(A)$, then $p(A)$ is invertible.
\begin{proof}
By Corollary \ref{cor:polynomial-spectral-mapping-matrix}, $0\notin\operatorname{spec}(p(A))$. By Corollary \ref{cor:matrix-invertible-iff-zero-not-spectrum}, $p(A)$ is invertible.\\
\hlt{Structural Commentary}: Invertibility of $p(A)$ is checked on spectral values of $A$.\\
\hlt{Key Insight}: Avoiding zero on the spectrum prevents $p(A)$ from having a kernel.
\end{proof}
\end{corollary}
